\subsection{ns3}
ns-3 is an open-source, discrete-event network simulator developed in C++ with support for integration with Python \cite{ns3}.
It is an end-to-end network simulator, alternative to OMNeT++, and is widely used in academic for research of communication networks.
ns-3 is developed to reflect the reality as close as possible with core concepts and abstractions that map closely to how the real communication network are built.

In ns-3, a \texttt{Node} is a fundamental network entity. It serves as a container for \texttt{Application}, \texttt{Protocols}, and \texttt{Network Devices}.
Each node can run an \texttt{Application} which is a user program that generate packet flows.
Ns-3 has several built-in traffic generator for different applications such as video streaming, social media, gaming, VoIP, FTP, etc.

A \texttt{Network Device} is an entity that connected to a \texttt{Channel} which represent the transmission medium between the devices.
The broad community has continuously contributed to the network simulator, which now supports different channel models such as Ethernet, Wi-Fi, LTE, WiMax, etc. \cite{ns3}.
Each time a node sends a packet, the \texttt{Channel} calls a statistical or stochastic model, such as Friis, Log-Distance, or 3GPP \texttt{Propagation\allowbreak LossModel} and \texttt{Propagation\allowbreak DelayModel}, to calculate the path loss and transmission delay for that particular transmission \cite{ns3,tr38901}.

Moreover, ns-3 provide mobility model such as \texttt{ConstantPositionMobilityModel} and \texttt{RandomWalk2MobilityModel} to simulate the movement of devices.
However, the propagation models provided in ns-3 are limited to statistical models and support for 5G NR and beyond is not supported yet in main framework.
